\chapter{Selections}\label{ch:select}
\begin{goals}
\textbf{In this chapter} you select measurements on a net: one at a time, round a cluster, within a cone of
a direction or a band about a great circle. You add to a selection and take from it, see it everywhere the
data is drawn, and keep it as a dataset of its own.
\end{goals}

A selection answers \enquote{which of these measurements are those?}: the poles in the cluster to the
south-west, the planes within 20° of the regional foliation, the lineations near a fold axis. wuffle draws
the chosen rows violet on every net and highlights them in the table. You can then read their statistics,
make them a dataset, or hand them on to the console.

\section{The tool strip}\label{sec:select-tools}

Every net carries a strip of four tools at its top left:

\begin{description}
\item[\ui{pick}] a click selects the measurement under the pointer. This is the tool a net starts with.
\item[\ui{lasso}] drag round the measurements you want.
\item[\ui{cone}] press on a direction and drag outwards: everything within that angle of it.
\item[\ui{band}] drag along a great circle: everything within a band either side of it. Choosing it shows
a field, \ui{±~10~°}, for the band's half-width, from 1° to 45°.
\end{description}

Hover over a tool for a reminder of what it does. A gesture selects in every \emph{visible} layer on that
net, so hide a layer in the tree to leave it out of a selection.

\section{Picking}\label{sec:pick}

With \ui{pick}, click a pole or a line. wuffle takes the nearest measurement within a few pixels of the
pointer:

\begin{itemize}
\item a click replaces that dataset's selection with that one measurement (other datasets' selections stay);
\item \key{Shift}+click adds it to the selection;
\item \key{Ctrl}+click (\key{⌘}+click on a Mac) toggles it: in if it was out, out if it was in;
\item a click on empty space clears the selection.
\end{itemize}

While you move over the net, the status line reads the direction under the pointer and the measurement
nearest to it: \ui{cursor · pole of 117/54 · line 297→36 · tocher \#152 117/54}. The \code{\#} is the
row number in the table.

\begin{note}
With \ui{pick} chosen, a press and drag turns the net instead of selecting (\cref{ch:rotation}). The
other three tools use the drag for their shape, so you rotate with \ui{pick} and select with the rest.
\end{note}

\section{Shapes}\label{sec:shapes}

\shot{select-cone}{A cone of 20° about the pole to \code{040/50}: the dashed outline, the 18 poles inside
it violet, and the status line counting them.}

The \ui{lasso}, \ui{cone} and \ui{band} draw a shape and select what is inside it when you let go:

\begin{description}
\item[Lasso.] Hold the button down and draw round the measurements. Letting go closes the loop. The status
line reports it as a \ui{polygon}, for example \ui{polygon: 38 of 200 measurements}.
\item[Cone.] Press on the direction the cone is about (a cluster's centre, a fold axis) and drag away from
it. The status line shows the cone's half-angle as you drag (\ui{cone: 20.0° about the axis}); everything
within that angle is selected when you let go. A cone smaller than about 0.2° is ignored, so a stray click
selects nothing.
\item[Band.] Press on one point of the great circle you mean and drag to another: the band follows the
great circle through both. Set its half-width first, in the field beside the tools. A band of ±10° about
a girdle picks the measurements that belong to it.
\end{description}

The last shape stays on the net, dashed, so you can see what made the selection. It goes when you draw
another shape or press \key{Esc}. \key{Esc} clears only the outline; the selection stays.

\subsection{Adding and taking away}

Hold a key as you start a shape:

\begin{itemize}
\item no key: the shape's measurements \emph{replace} the selection;
\item \key{Shift}: they are \emph{added} to it;
\item \key{Alt} (\key{⌥} on a Mac): they are \emph{taken out} of it.
\end{itemize}

The status line at the foot of the window says which it was: \ui{band: 58 of 200 measurements added},
\ui{cone: 16 of 200 measurements removed}. That count is what the shape caught; to its left the window
keeps the total, \ui{1 dataset · 200 measurements · 56 selected}.

\shot{select-band}{The band tool chosen, with its half-width field. A band of ±10° about \code{040/50}
added to the cone of \cref{fig:select-cone} (\code{select + band 040/50 10}, the same as a \key{Shift} drag): 45 selected.}

\section{On the sphere, not the screen}\label{sec:select-sphere}

A shape is not drawn on the picture of the net; it is drawn on the sphere of directions. A cone is an
axis and an angle, a band a great circle and a half-width, and a lasso a spherical polygon whose edges are
great-circle arcs. wuffle tests each measurement against the shape in three dimensions.

\begin{why}
A shape on the screen would select different measurements on an equal-angle net than on an equal-area
one, and different ones again once the net is turned. A shape on the sphere means the same measurements
whatever the projection, the hemisphere or the rotation, so a selection made on a rotated net is the one you
would have made on the plain net (\cref{sec:rotation-true}).
\end{why}

Every test is \emph{axial}: a pole and the opposite end of the same line are one measurement. A cone about
a shallow direction near the edge of the net catches the poles just inside the opposite edge too, which
is why its outline can show as two arcs on opposite sides of the primitive. A lasso has to fit inside a
hemisphere, which any lasso drawn on one net does.

\section{The selection belongs to the data}\label{sec:select-data}

A selection is a set of \emph{rows} of a dataset, not marks on a net. So:

\begin{itemize}
\item the same rows are violet on every net that draws the dataset, and every layer that shows them;
\item the table highlights them, and selecting rows in the table selects them on the nets;
\item it is saved with the project (\file{selection/<name>.rows} in the project file, \cref{ch:projects}).
\end{itemize}

Choosing rows does not change the project, so it is never an undo step of its own. Undoing a change puts
the selection back as it was before that change (\cref{sec:undo}).

\begin{note}
The statistics card is for the whole dataset, not the selection: selecting rows leaves its numbers as
they were. To see the statistics of a selection, make it a dataset (\cref{sec:select-dataset}).
\end{note}

\section{A selection as a dataset}\label{sec:select-dataset}

Right-click a layer in the tree and choose \ui{Selection → new dataset} (greyed out while nothing of that
dataset is selected). The selected rows become a new dataset, \code{tocher · selection}, with the same
columns, kind and reading of the attitude. It appears in the tree's \ui{data} section with an arrow, and
it remembers where it came from: its lineage names \code{tocher} and the rows it took. It is one undo step,
\ui{Undo selection → dataset}.

Draw it on a net and its statistics card appears beside the full set's (\cref{ch:stats}). This is how you
split a dataset into domains: select the measurements of each, make each a dataset, and compare.

\section{From the console}\label{sec:select-console}

The console's \code{select} draws the same shapes from typed attitudes (\cref{sec:commands}), on every
visible layer of the active net, and draws the outline too:

\begin{lstlisting}
select cone 040/50 20       # within 20 degrees of the pole to 040/50
select + band 040/50 10     # add those within 10 degrees of the plane 040/50
select - cone 040/50 10     # take out the inner 10 degrees
select invert               # everything that wasn't, and nothing that was
select none
\end{lstlisting}

A plane's cone is about its pole; a line's (\code{'150>60'}, quoted, or \code{150→60}) about the line.
Each command answers with what its shape caught: \ui{cone: 18 selected}.

The selection is also a file. \file{selection/tocher.rows} lists the selected row numbers, one per line,
counting from 1; write numbers to it to select those rows, \code{rm} it to clear
(\cref{sec:project-write}). \file{selection/tocher.csv} holds the selected rows themselves.

\begin{tryit}
\begin{enumerate}
\item Open the Tocher sample and type \code{select cone 040/50 20} in the console. 18 poles turn violet,
inside a dashed circle.
\item Choose \ui{band} on the net, check the field says \ui{10}, then hold \key{Shift} and drag from
the south-west edge of the net towards the north-west. A band joins the selection and the status bar's
count grows.
\item Click the \ui{Net 1} tab (so the keys aren't going to the console) and press \key{Esc}. The
outlines go; the selection stays.
\item Right-click \ui{tocher} in the tree and choose \ui{Selection → new dataset}. The table shows the new
dataset, the rows you chose and nothing else.
\item Type \code{cat selection/tocher.rows} in the console to see their row numbers.
\end{enumerate}
\end{tryit}
